% Test fixture: heavy math paper
\documentclass{article}
\usepackage{amsmath,amssymb,amsthm}
\usepackage{mathrsfs}
\usepackage{bm}

\begin{document}

\title{Stress Test: Complex Mathematical Typesetting}
\author{Test Author}
\maketitle

\section{Operator Stacking}

\begin{equation}
    \sum_{\substack{i=1 \\ i \neq j}}^{n} \prod_{k=1}^{m}
    \frac{\partial^2 f}{\partial x_i \partial x_k} \cdot
    \oint_{\gamma} \mathscr{F}(\zeta) \, d\zeta
\end{equation}

\section{Nested Fractions}

\begin{equation}
    \cfrac{1}{1 + \cfrac{1}{1 + \cfrac{1}{1 + \cfrac{1}{x}}}}
\end{equation}

\section{Matrix Operations}

\begin{equation}
    \det \begin{vmatrix}
        a_{11} & a_{12} & a_{13} \\
        a_{21} & a_{22} & a_{23} \\
        a_{31} & a_{32} & a_{33}
    \end{vmatrix} =
    \sum_{\sigma \in S_3} \text{sgn}(\sigma)
    \prod_{i=1}^{3} a_{i,\sigma(i)}
\end{equation}

\section{Commutative Diagrams (as aligned)}

\begin{equation}
    \begin{aligned}
        A &\xrightarrow{f} B \\
        {\scriptstyle g}\downarrow &\quad \downarrow{\scriptstyle h} \\
        C &\xrightarrow{k} D
    \end{aligned}
\end{equation}

\section{Multi-line Derivations}

\begin{align}
    \nabla \times (\nabla \times \mathbf{E})
        &= \nabla(\nabla \cdot \mathbf{E}) - \nabla^2 \mathbf{E} \\
        &= -\frac{\partial}{\partial t}(\nabla \times \mathbf{B}) \\
        &= -\mu_0 \epsilon_0 \frac{\partial^2 \mathbf{E}}{\partial t^2} \\
    \Rightarrow \quad \nabla^2 \mathbf{E}
        &= \mu_0 \epsilon_0 \frac{\partial^2 \mathbf{E}}{\partial t^2}
\end{align}

\section{Accents and Decorations}

\begin{equation}
    \hat{x}, \bar{x}, \vec{x}, \dot{x}, \ddot{x},
    \tilde{x}, \widehat{xyz}, \widetilde{abc},
    \overbrace{a+b+c}^{\text{sum}},
    \underbrace{x \cdot y \cdot z}_{\text{product}}
\end{equation}

\section{Blackboard Bold and Script}

\begin{equation}
    \mathbb{R}^n, \quad \mathbb{C}^{m \times n}, \quad
    \mathbb{Z}_p, \quad \mathbb{Q}(\sqrt{2}), \quad
    \mathscr{L}(V, W), \quad \mathcal{O}(n \log n), \quad
    \mathfrak{g} = \mathfrak{sl}(2, \mathbb{C})
\end{equation}

\section{Cases and Piecewise}

\begin{equation}
    |x| = \begin{cases}
        x & \text{if } x \geq 0 \\
        -x & \text{if } x < 0
    \end{cases}
\end{equation}

\begin{equation}
    \binom{n}{k} = \frac{n!}{k!(n-k)!} \quad \text{for } 0 \leq k \leq n
\end{equation}

\end{document}
